%; whizzy paragraph -pdf xpdf -latex ./whizzypdfptex.sh
%; whizzy-paragraph "^\\\\begin{frame}"
% latex beamer presentation.
% platex, latex-beamer でコンパイルすることを想定。 

%     Tokyo Debian Meeting resources
%     Copyright (C) 2009 Junichi Uekawa
%     Copyright (C) 2009 Nobuhiro Iwamatsu

%     This program is free software; you can redistribute it and/or modify
%     it under the terms of the GNU General Public License as published by
%     the Free Software Foundation; either version 2 of the License, or
%     (at your option) any later version.

%     This program is distributed in the hope that it will be useful,
%     but WITHOUT ANY WARRANTY; without even the implied warreanty of
%     MERCHANTABILITY or FITNESS FOR A PARTICULAR PURPOSE.  See the
%     GNU General Public License for more details.

%     You should have received a copy of the GNU General Public License
%     along with this program; if not, write to the Free Software
%     Foundation, Inc., 51 Franklin St, Fifth Floor, Boston, MA  02110-1301 USA

\documentclass[cjk,dvipdfmx,12pt]{beamer}
\usetheme{Tokyo}
\usepackage{monthlypresentation}

%  preview (shell-command (concat "evince " (replace-regexp-in-string
%  "tex$" "pdf"(buffer-file-name)) "&")) 
%  presentation (shell-command (concat "xpdf -fullscreen " (replace-regexp-in-string "tex$" "pdf"(buffer-file-name)) "&"))
%  presentation (shell-command (concat "evince " (replace-regexp-in-string "tex$" "pdf"(buffer-file-name)) "&"))

%http://www.naney.org/diki/dk/hyperref.html
%日本語EUC系環境の時
\AtBeginDvi{\special{pdf:tounicode EUC-UCS2}}
%シフトJIS系環境の時
%\AtBeginDvi{\special{pdf:tounicode 90ms-RKSJ-UCS2}}

\title{第144回東京エリアDebian勉強会 \\　\\preseedでDebianを自動インストールをしてみよう}
\subtitle{}
\author{Norimitsu Sugimoto (杉本 典充) \\dictoss@live.jp}
\date{2016-10-15}
\logo{\includegraphics[width=8cm]{image200607/openlogo-light.eps}}

\begin{document}

\frame{\titlepage{}}

%\section{}

\begin{frame}{アジェンダ}
  \begin{itemize}
  \item 自己紹介
  \item debianをインストールする方法
  \item pressedとは
  \item preseedを使ってdebianをインストールする
  \item netinstイメージを使った自動インストール
  \item virt-installを使った自動インストール
  \item まとめ
  \item 参考資料
  \end{itemize}
\end{frame}


\begin{frame}{自己紹介}
  \begin{itemize}
  \item Norimitsu Sugimoto (杉本 典充)
  \item dictoss@live.jp
  \item Twitter: @dictoss
  \item Debian使って10年以上、FreeBSD使って5年以上
  \item Debian GNU/kFreeBSDが気になっておりウォッチ中
  \item tilegxに少し手を出したがなかなか手ごわい
  \item 仕事はソフトウェア開発者をやってます
  \item pythonとDjangoの組み合わせで使うことが多いです
  \end{itemize}
\end{frame}


\emtext{debianをインストールする方法}

\begin{frame}[containsverbatim]{debianをインストールする方法}
  \begin{itemize}
  \item Debian Installerを使うことが多い
    \begin{itemize}
    \item 全パッケージを収録したISOイメージ群。CD、DVD、Blu-ray形式
    \item netinstイメージ。CDのISO形式、USBメモリ形式
    \item vmlinuz、initrd.gzの単体ファイル。ミラーサイトに置いてあります
    \end{itemize}
  \item 日常利用するPCへインストールするならnetinstイメージの利用がおすすめ
  \end{itemize}
\end{frame}


\emtext{preseed}

\begin{frame}[containsverbatim]{preseedとは}
  \begin{itemize}
  \item Debian Installerの「インストールの実行中に手動で回答を入力せずに、インストールプロセス中の質問の答を設定する方法を提供」する機能
  \item \url{https://www.debian.org/releases/jessie/amd64/apb.html.ja}
  \item preseed.cfg = インストール処理のパラメータを定義したファイル
  \end{itemize}
\end{frame}

\begin{frame}[containsverbatim]{preseedの区分}
  \begin{itemize}
  \item initrd preseed
    \begin{itemize}
    \item Debian Installerのinitrd.gzの中に/preseed.cfgを配置する
    \item インストーラ起動直後の選択肢から/preseed.cfgの定義に従い自動インストールする
    \end{itemize}
  \end{itemize}
\end{frame}

\begin{frame}[containsverbatim]{preseedの区分}
  \begin{itemize}
  \item file preseed
  \begin{itemize}
    \item Debian InstallerのISOファイルの/preseed.cfgを配置する（＝要リマスタリング）
    \item Debian InstallerのUSBメモリの/preseed.cfgを配置する
    \item kernelブートパラメータでファイルの場所を指定する（preseed/file、preseed/file/checksum）
    \end{itemize}
  \end{itemize}
\end{frame}

\begin{frame}[containsverbatim]{preseedの区分}
  \begin{itemize}
  \item network preseed
    \begin{itemize}
    \item IPアドレスの取得または設定後にpreseed.cfgをwgetでダウンロードする
    \item kernelブートパラメータでファイルの場所を指定する（preseed/url、preseed/url/checksum）
    \item IPアドレスを取得する前の選択肢を出さないようにするため追加のkernelブートパラメータも指定が必要
    \end{itemize}
  \end{itemize}
\end{frame}

\begin{frame}[containsverbatim]{preseedファイルの仕様}
  \begin{itemize}
  \item 仕様はwebサイトに説明があります
    \begin{itemize}
    \item 「B.4. 事前設定ファイルの内容 (jessie 用)」
    \item \url{https://www.debian.org/releases/stable/s390x/apbs04.html.ja}
    \item 「B.2.4. preseed で利用できるエイリアス」
    \item \url{https://www.debian.org/releases/stable/s390x/apbs02.html.ja#preseed-aliases}
    \end{itemize}
  \item Debian 8(Jessie)用のサンプルがあります
    \begin{itemize}
    \item \url{https://www.debian.org/releases/jessie/example-preseed.txt}
    \end{itemize}
  \end{itemize}
\end{frame}


\begin{frame}[containsverbatim]{インストール中のカスタムコマンド実行}
  \begin{itemize}
  \item preseed/early\_command
    \begin{itemize}
    \item 全体のインストール処理が始まる前にコマンドを実行する
    \end{itemize}
  \item 処理名/early\_command　
    \begin{itemize}
      \item ある処理が始まる前にコマンドを実行する
      \item 例） partman/early\_command 
    \end{itemize}
  \item preseed/late\_command
    \begin{itemize}
    \item インストールが完了し、再起動する前にコマンドを実行する
    \item インストール先ディスクの/は、/targetとしてマウント中
    \end{itemize}
  \end{itemize}
\end{frame}


\emtext{preseedを使ってdebianをインストールする}

\begin{frame}[containsverbatim]{インストールテスト環境}
  \begin{itemize}
  \item Debian GNU/Linux 8のKVMホスト環境とした仮想マシンへインストールする
  \item 使用するpreseed.cfgは資料参照
  \end{itemize}
\end{frame}


\begin{frame}[containsverbatim]{netinstでpreseed}
  \begin{itemize}
  \item netinstのISOイメージなどからブートした場合に使う
  \item network preseedとして動作する
  \item Debian Installerのkernelブートパラメータに以下を入力して起動
  \end{itemize}
  \begin{commandline}
    auto=true priority=critical \
    url=http://192.168.122.41/preseed.cfg \
    preseed/url/checksum=cdd4bc6c41bd8d96dbcd37c6df442222
  \end{commandline}
  \begin{itemize}
  \item preseed/url/checksumはmd5sumです（必須）
  \item md5sumが誤っている場合、preseedファイルは使用されない
  \end{itemize}
\end{frame}


\begin{frame}[containsverbatim]{virt-installでpreseed}
  \begin{itemize}
  \item virt-installの機能にあるinitrd-injectオプションを使う
  \item initrd preseedとして動作する
  \item ファイル名は必ずpreseed.cfgにすること
  \end{itemize}
  \begin{commandline}
$ sudo virt-install \
    --name deb-preseed-1 \
    (省略) \
    --os-type linux --os-variant generic \
    --console pty,target_type=serial \
    --location 'http://ftp.jp.debian.org/debian/dists/jessie/main/installer-amd64/' \
    --initrd-inject=/var/lib/libvirt/images/preseed.cfg \
    --extra-args 'console=ttyS0,115200n8 serial'
  \end{commandline}
\end{frame}


\emtext{まとめ}

\begin{frame}[containsverbatim]{まとめ}
  \begin{itemize}
  \item preseed機能を試してみました
  \item 仮想化環境へ日常的にDebianをインストールしている場合は効率が上がると思います
  \item うまく活用してdebianライフをお楽しみください
  \end{itemize}
\end{frame}

  
\begin{frame}{参考情報}
  \begin{itemize}
  \item 「DebianInstallerPreseed」 \url{https://wiki.debian.org/DebianInstaller/Preseed}
  \item 「付録B preseed を利用したインストールの自動化」 \url{https://www.debian.org/releases/jessie/amd64/apb.html.ja}
  \end{itemize}
\end{frame}

\end{document}


;;; Local Variables: ***
;;; outline-regexp: "\\([ 	]*\\\\\\(documentstyle\\|documentclass\\|emtext\\|section\\|begin{frame}\\)\\*?[ 	]*[[{]\\|[]+\\)" ***
;;; End: ***
